\documentclass[12pt]{article}
\makeatletter
\def\input@path{{../../style/}}
\makeatother

\usepackage{../../style/psetconfig}
\title{Homework 4}
\author{Songyu Ye}
\date{\today}

\begin{document}
\psettitle

\begin{problem}[1]
Prove that the map
\[
S^{4n-1}\to S^{2n}
\]
defined by the Whitehead square of the generator has Hopf invariant \(2\).
\end{problem}

\begin{solution}
    The mapping cone of the Whitehead bracket $S^{2n-1} \to S^n$ is $C_{[\iota,\iota]}=S^n\cup_{[\iota,\iota]} e^{2n}$ the space obtained by folding $S^n \vee S^n \subset S^n \times S^n$ to a single $S^n$. We have
\[
H^n(S^n\times S^n) = \mathbb{Z}a \oplus \mathbb{Z}b,
\]
and
\[
H^{2n}(S^n\times S^n) = \mathbb{Z}\,ab.
\]
Observe that the class of the folded sphere \(x \in H^n(C_{[\iota,\iota]},\mathbb{Z}) \cong \Z\) pulls back to \(a+b\) in \(S^n\times S^n\),
and consequently \(x^2\) goes to \((a+b)^2\). Now,
\[
(a+b)^2 = a^2 + b^2 + ab + ba = ab + ba = \bigl(1 + (-1)^n\bigr)ab.
\]

In particular, the Hopf invariant is
\[
H([\iota,\iota]) =
\begin{cases}
2 & \text{if \(n\) is even},\\
0 & \text{if \(n\) is odd}.
\end{cases}
\]
\end{solution}

\begin{problem}[2 (Hatcher AT, 4.2.34)]
Let \(p:S^3\to S^2\) be the Hopf bundle and \(q:T^3\to S^3\) be the quotient map collapsing the complement of a ball in the \(3\)-dimensional torus \(T^3=S^1\times S^1\times S^1\) to a point. Show that \(pq:T^3\to S^2\) induces the zero map on homotopy and reduced homology, but is \emph{not} homotopic to a constant map.
\end{problem}

\begin{solution}
Let \(p:S^3\to S^2\) be the Hopf fibration and let
\[
q:T^3\to S^3
\]
be the quotient collapsing the complement of an embedded \(3\)-ball in \(T^3=S^1\times S^1\times S^1\) to a point. Set \(f:=p\circ q\).

The cellular approximation theorem allows us to homotope \(q\) to a cellular map, which we also denote by \(q\). This $q$ must collapse the \(2\)-skeleton of \(T^3\) a point and so $q_*=0$ for $i<3$. Also $q_*$ is an isomorphism in degree $3$ since it sends the fundamental class of the chosen ball to the fundamental class of $S^3$. On the other hand, $p_*$ is $0$ in degree $3$ since $H_3(S^2)=0$. Therefore, the composition \(f=p\circ q\) induces the zero map on reduced homology.

The universal cover of $T^3$ is $\R^3$, which is contractible, so $\pi_i(T^3)=0$ for $i>1$. Also $\pi_1(T^3)=\Z^3$ and $\pi_1(S^2) = 0$ so $f_*$ is the zero map on all homotopy groups.

Suppose for contradiction that \(f\) is homotopic to a constant map. Then there exists a homotopy
\[
H:T^3\times I\to S^2
\]
with \(H(-,0)=f\) and \(H(-,1)=c\) constant.

Since \(p:S^3\to S^2\) is a fibration and \(q\) is a lift of \(f\), the homotopy lifting property produces a lift
\[
\widetilde H:T^3\times I\to S^3
\]
such that
\[
\widetilde H(-,0)=q
\qquad\text{and}\qquad
p\circ \widetilde H = H.
\]

At time \(t=1\), since \(H(-,1)=c\) is constant, the map \(\widetilde H(-,1)\) must land in a single fiber of \(p\), which is a circle \(S^1\subset S^3\). Thus \(\widetilde H(-,1)\) factors as
\[
T^3 \longrightarrow S^1 \hookrightarrow S^3.
\]

Therefore \(q\) is homotopic (as a map into \(S^3\)) to a map factoring through \(S^1\). But this is impossible: the map \(q\) has degree \(1\), so
\[
q_*:H_3(T^3)\to H_3(S^3)
\]
is an isomorphism, whereas any map \(T^3\to S^1\subset S^3\) induces the zero map on \(H_3\) since \(H_3(S^1)=0\).
\end{solution}

\begin{problem}[3]
For a relatively finite CW pair \((X,A)\) (meaning that \(X\) is built by attaching finitely many cells to \(A\)) and any space \(Y\), show that the restriction map
\[
\operatorname{Map}(X,Y)\to \operatorname{Map}(A,Y)
\]
is a fibration.
\end{problem}

\begin{solution}
Let
\[
r:\operatorname{Map}(X,Y)\to \operatorname{Map}(A,Y)
\]
be the restriction map. We show that \(r\) is a fibration.

It suffices to verify the homotopy lifting property with respect to disks. So suppose we are given a commutative diagram
\[
\begin{array}{ccc}
D^n & \xrightarrow{\,g\,} & \operatorname{Map}(X,Y) \\
\big\downarrow && \big\downarrow r \\
D^n\times I & \xrightarrow{\,H\,} & \operatorname{Map}(A,Y)
\end{array}
\]
where \(H(u,0)=r(g(u))\) for all \(u\in D^n\). We must construct a lift
\[
\widetilde H:D^n\times I\to \operatorname{Map}(X,Y)
\]
such that \(\widetilde H(u,0)=g(u)\) and \(r\circ \widetilde H=H\).

By adjunction, \(g\) corresponds to a map
\[
\widehat g:D^n\times X\to Y,
\]
and \(H\) corresponds to a map
\[
\widehat H:D^n\times I\times A\to Y.
\]
The condition \(H(u,0)=r(g(u))\) says exactly that
\[
\widehat H(u,0,a)=\widehat g(u,a)
\qquad
\text{for all }u\in D^n,\ a\in A.
\]

Hence \(\widehat g\) and \(\widehat H\) together define a map
\[
F:(D^n\times X\times\{0\})\cup(D^n\times A\times I)\to Y.
\]
So it remains to extend \(F\) to all of \(D^n\times X\times I\).

Now \((X,A)\) is a relatively finite CW pair, so \(X\) is obtained from \(A\) by attaching finitely many cells. Therefore \((D^n\times X,D^n\times A)\) is again a relative CW pair, hence the inclusion
\[
D^n\times A\hookrightarrow D^n\times X
\]
is a cofibration. It follows that the inclusion
\[
(D^n\times X\times\{0\})\cup(D^n\times A\times I)
\hookrightarrow
D^n\times X\times I
\]
is also a cofibration, so the homotopy extension property applies. Thus \(F\) extends to a map
\[
\widetilde F:D^n\times X\times I\to Y.
\]

Adjointing back, \(\widetilde F\) corresponds to a map
\[
\widetilde H:D^n\times I\to \operatorname{Map}(X,Y)
\]
with
\[
\widetilde H(u,0)=g(u)
\qquad\text{and}\qquad
r(\widetilde H(u,t))=H(u,t).
\]
Thus \(r\) has the homotopy lifting property for disks, hence is a fibration.
\end{solution}
\begin{problem}[4]
Given fibrations \(E,E'\to B\), show that any fiber-preserving map
\[
f:E\to E'
\]
which is a homotopy equivalence (on total spaces) is a \emph{fiber homotopy equivalence}, meaning there exists a fiber-preserving map
\[
g:E'\to E
\]
such that \(g\circ f\) and \(f\circ g\) are both homotopic to the identities via fiber-preserving maps.
\end{problem}

\begin{solution}
We first record a standard lemma.

\medskip

\noindent\textbf{Lemma.}
Let \(p:E\to B\) be a fibration, and let \(u:X\to B\) be a map. If
\[
f_0,f_1:X\to E
\]
satisfy
\[
p\circ f_0=p\circ f_1=u,
\]
and if \(f_0\simeq f_1\) as ordinary maps \(X\to E\), then in fact \(f_0\) and \(f_1\) are homotopic through fiber-preserving maps, i.e. there exists
\[
\widetilde H:X\times I\to E
\]
with
\[
\widetilde H(-,0)=f_0,\qquad \widetilde H(-,1)=f_1,
\qquad p\widetilde H(x,t)=u(x).
\]

\begin{proof}
Let
\[
H:X\times I\to E
\]
be a homotopy from \(f_0\) to \(f_1\). Its projection
\[
pH:X\times I\to B
\]
is a homotopy from \(u\) to \(u\). Define
\[
G:X\times I\times I\to B,\qquad G(x,s,t)=pH(x,(1-s)t).
\]
Then
\[
G(x,0,t)=pH(x,t),\qquad G(x,1,t)=u(x).
\]
So \(G\) is a homotopy, relative \(X\times\{0,1\}\), from \(pH\) to the constant homotopy at \(u\).

Now define
\[
\Phi:X\times I\times\{0\}\;\cup\;X\times\{0\}\times I\;\cup\;X\times\{1\}\times I \to E
\]
by
\[
\Phi(x,s,0)=f_0(x),\qquad
\Phi(x,0,t)=H(x,t),\qquad
\Phi(x,1,t)=f_1(x).
\]
These formulas agree on overlaps, and one checks that $p\Phi=G$
on this subspace. Since
\[
(X\times I\times\{0\})\cup(X\times\{0,1\}\times I)\hookrightarrow X\times I\times I
\]
is a cofibration and \(p\) is a fibration, the homotopy lifting property yields an extension
\[
\widetilde\Phi:X\times I\times I\to E
\]
with
$p\widetilde\Phi=G$.
Set $\widetilde H(x,t):=\widetilde\Phi(x,1,t)$.
Then $\widetilde H(-,0)=f_0$, $\widetilde H(-,1)=f_1$
and \[p\widetilde H(x,t)=G(x,1,t)=u(x)\]
for all \((x,t)\). Thus \(\widetilde H\) is a fiber-preserving homotopy.
\end{proof}

\medskip

Now let
\[
p:E\to B,\qquad p':E'\to B
\]
be fibrations, and let
\[
f:E\to E'
\]
be fiber-preserving, so \(p'f=p\). Assume \(f\) is a homotopy equivalence of total spaces. Choose an ordinary homotopy inverse
\[
h:E'\to E
\]
with
\[
hf\simeq \mathrm{id}_E,\qquad fh\simeq \mathrm{id}_{E'}.
\]

We first deform \(h\) to a fiber-preserving map. Since \(p'f\simeq p' \circ \mathrm{id}_{E'}\) and \(fh\simeq \mathrm{id}_{E'}\), we have
\[
ph = p'fh \simeq p'.
\]
Apply the homotopy lifting property to the homotopy \(ph\simeq p'\), starting with the map \(h\). This gives a homotopy from \(h\) to a map
\[
g:E'\to E
\]
such that
\[
pg=p'.
\]
So \(g\) is fiber-preserving.

Now \(g\simeq h\), hence
\[
gf\simeq hf\simeq \mathrm{id}_E
\]
and
\[
fg\simeq fh\simeq \mathrm{id}_{E'}.
\]
Moreover both pairs of maps are fiber-preserving:
\[
p(gf)=p'f=p=p(\mathrm{id}_E),
\]
\[
p'(fg)=p'g=p'=p'(\mathrm{id}_{E'}).
\]
Therefore, by the lemma, these ordinary homotopies can be replaced by fiber-preserving homotopies. Thus
\[
gf\simeq_B \mathrm{id}_E,\qquad fg\simeq_B \mathrm{id}_{E'}.
\]

Hence \(g\) is a fiber-preserving homotopy inverse to \(f\), so \(f\) is a fiber homotopy equivalence.
\end{solution}


\begin{problem}[5]
Show that a fibration over a \emph{contractible} base is fiber homotopy equivalent to a product space. Show the same for a fibration with contractible fiber with CW base. Assuming that a fiber bundle over a paracompact base is a fibration (for example when the fibers are smooth manifolds), conclude that every fiber bundle over a contractible paracompact base is a product bundle.
\end{problem}

\begin{solution}
First suppose \(p:E\to B\) is a fibration and \(B\) is contractible. Choose a basepoint \(b_0\in B\), let
\[
F=p^{-1}(b_0),
\]
and let
\[
H:B\times I\to B
\]
be a contraction of \(B\) to \(b_0\), so \(H(b,0)=b\) and \(H(b,1)=b_0\).

Now \(H\) is a homotopy from \(\mathrm{id}_B\) to the constant map \(c:B\to B\), \(c(b)=b_0\). Pulling back the fibration \(p\) along these two maps gives:
\[
(\mathrm{id}_B)^*E \cong E,
\qquad
c^*E \cong B\times F.
\]
Since pullbacks of a fibration along homotopic maps are fiber homotopy equivalent, it follows that
\[
E \simeq_B B\times F.
\]
Thus a fibration over a contractible base is fiber homotopy equivalent to a product.

Now suppose instead that \(p:E\to B\) has contractible fiber \(F\). We claim that \(p\) is a fiber homotopy equivalence with inverse given by a section \(s:B\to E\). Indeed, over a CW decomposition of \(B\), the obstruction to extending a section over an \(n\)-cell lies in \(\pi_{n-1}(F)\), and these groups vanish since \(F\) is contractible. Hence a section \(s\) exists. Then
\[
p\circ s=\mathrm{id}_B.
\]
Also \(s\circ p:E\to E\) is fiber-preserving, and on each fiber \(E_b\) it is the constant map to \(s(b)\). Since each fiber is contractible, \(s\circ p\) is fiberwise homotopic to the identity. Therefore
\[
s\circ p \simeq_B \mathrm{id}_E,
\]
so \(p\) is a fiber homotopy equivalence. Hence \(E\) is fiber homotopy equivalent to \(B\).


Finally, let \(p:E\to B\) be a fiber bundle over a contractible paracompact base \(B\), and choose \(b_0\in B\). Let \(c:B\to B\) be the constant map \(c(b)=b_0\). Since \(B\) is contractible, \(\mathrm{id}_B\) is homotopic to \(c\). By the bundle homotopy theorem for bundles over paracompact bases, pullbacks along homotopic maps are isomorphic. Hence
\[
E \cong \mathrm{id}_B^*E \cong c^*E.
\]
If \(F=p^{-1}(b_0)\), then \(c^*E\cong B\times F\). Therefore
\[
E\cong B\times F
\]
as bundles over \(B\). Thus every fiber bundle over a contractible paracompact base is a product bundle.

\end{solution}

\begin{problem}[6 (Hatcher AT, 4.2.17)]
Show that the map
\[
[X,Y]\to \operatorname{Hom}(\pi_n X,\pi_n Y)
\]
sending \([f]\mapsto f_*\) is a bijection, if \(X\) is an \((n-1)\)-connected CW complex and \(Y\) is a path-connected space with \(\pi_i Y=0\) for \(i>n\). Deduce that CW models for \(K(G,n)\) are unique up to homotopy equivalence.
\end{problem}

\begin{solution}
Since \(X\) is \((n-1)\)-connected, we may choose a CW structure on \(X\) with
\[
X^{n-1}=*.
\]
Then \(X^n\) is a wedge of \(n\)-spheres,
\[
X^n \cong \bigvee_\alpha S^n_\alpha,
\]
and the quotient \(X^n/X^{n-1}\) is the same wedge.

A based map
\[
f_n:X^n\to Y
\]
is therefore determined up to homotopy by its restrictions to the \(n\)-cells, hence by a cellular \(n\)-cochain with values in \(\pi_n(Y)\). More concretely, such a map determines a homomorphism
\[
C_n(X)\to \pi_n(Y),
\]
and two such maps are homotopic rel \(X^{n-1}\) exactly when they induce the same homomorphism
\[
H_n(X)\to \pi_n(Y).
\]
Thus homotopy classes of maps \(X^n\to Y\) correspond naturally to
\[
\operatorname{Hom}(H_n(X),\pi_n(Y)).
\]

Now we claim that every map \(X^n\to Y\) extends over all higher skeleta of \(X\), and that any two such extensions are homotopic. Indeed, the obstruction to extending from \(X^k\) to \(X^{k+1}\) lies in
\[
H^{k+1}(X;\pi_k(Y)).
\]
But for \(k\ge n\), we have \(\pi_k(Y)=0\) unless \(k=n\), and for \(k=n\) the map is already defined on \(X^n\). Hence all higher obstructions vanish. Similarly, the obstruction groups controlling uniqueness of extension involve \(\pi_k(Y)\) for \(k>n\), so they also vanish. Therefore restriction induces a bijection
\[
[X,Y]\xrightarrow{\;\cong\;} [X^n,Y].
\]

Combining this with the previous paragraph gives
\[
[X,Y]\cong \operatorname{Hom}(H_n(X),\pi_n(Y)).
\]
Since \(X\) is \((n-1)\)-connected, the Hurewicz theorem gives
\[
\pi_n(X)\cong H_n(X),
\]
so we obtain a natural bijection
\[
[X,Y]\cong \operatorname{Hom}(\pi_n(X),\pi_n(Y)).
\]

It remains to note that this bijection is exactly the map \([f]\mapsto f_*\). Indeed, if \(f:X\to Y\), then its restriction to \(X^n\) determines the induced homomorphism
\[
H_n(X)\to \pi_n(Y),
\]
and under the Hurewicz isomorphism \(H_n(X)\cong \pi_n(X)\), this is precisely the induced map
\[
f_*:\pi_n(X)\to \pi_n(Y).
\]
Thus the stated map is a bijection.

\medskip

Now let \(X\) and \(Y\) be CW models for \(K(G,n)\). Then \(X\) is \((n-1)\)-connected and \(Y\) is path-connected with \(\pi_i(Y)=0\) for \(i>n\), so the first part applies:
\[
[X,Y]\cong \operatorname{Hom}(\pi_n(X),\pi_n(Y)).
\]
Since
\[
\pi_n(X)\cong G\cong \pi_n(Y),
\]
the identity map on \(G\) is induced by some map
\[
f:X\to Y.
\]
Similarly there is a map
\[
g:Y\to X
\]
inducing the identity on \(\pi_n\). Then
\[
(g\circ f)_*=\mathrm{id}_{\pi_n(X)},
\qquad
(f\circ g)_*=\mathrm{id}_{\pi_n(Y)}.
\]
Applying the bijection again, we conclude that
\[
g\circ f\simeq \mathrm{id}_X,
\qquad
f\circ g\simeq \mathrm{id}_Y.
\]
Hence \(f\) and \(g\) are homotopy inverses. Therefore CW models for \(K(G,n)\) are unique up to homotopy equivalence.
\end{solution}

\begin{problem}[7 (var.\ of Hatcher AT, 4.3.6 \& 4.3.7)]
For an Abelian group \(G\), Question 6 gives a map
\[
\mu:K(G,n)\times K(G,n)\to K(G,n)
\]
corresponding to the addition on \(G\). Show that \(\mu\) defines an \(H\)-space structure on \(K(G,n)\) that is commutative and associative up to homotopy.

For any CW space \(X\), use \(\mu\) to define an addition on \([X,K(G,n)]\) by
\[
[f]+[g]:=\mu(f\times g)
\]
and show that under the bijection
\[
[X,K(G,n)]\cong H^n(X;G)
\]
this corresponds to the usual addition in cohomology.
\end{problem}

\begin{solution} Let \(Y=K(G,n)\). Let \(0\in G\) denote the identity element, and let \(*\in Y\) be the basepoint. The constant map
\[
e:* \to Y
\]
corresponds to \(0\in \pi_n(Y)=G\). Consider the maps
\[
\mu\circ(\mathrm{id},e),\qquad \mu\circ(e,\mathrm{id}):Y\to Y.
\]
On \(\pi_n(Y)\cong G\), both induce the homomorphism
\[
a\longmapsto a+0=a.
\]
By Question 6, these maps are homotopic to \(\mathrm{id}_Y\). Thus \(*\) is a homotopy identity for \(\mu\), so \((Y,\mu)\) is an \(H\)-space. To show commutativity, let
\[
\tau:Y\times Y\to Y\times Y,\qquad \tau(x,y)=(y,x).
\]
Then the two maps
\[
\mu,\qquad \mu\circ\tau : Y\times Y\to Y
\]
induce on \(\pi_n(Y\times Y)\cong G\oplus G\) the homomorphisms
\[
(a,b)\mapsto a+b,\qquad (a,b)\mapsto b+a.
\]
Since \(G\) is Abelian, these are equal. Hence by Question 6,
\[
\mu\circ\tau \simeq \mu.
\]
So \(\mu\) is commutative up to homotopy. To show associativity, consider the two maps
\[
\mu\circ(\mu\times \mathrm{id}),\qquad \mu\circ(\mathrm{id}\times \mu):
Y\times Y\times Y\to Y.
\]
Since
\[
\pi_n(Y\times Y\times Y)\cong G\oplus G\oplus G,
\]
the first induces
\[
(a,b,c)\longmapsto (a+b)+c,
\]
while the second induces
\[
(a,b,c)\longmapsto a+(b+c).
\]
These are equal because addition in \(G\) is associative. Therefore, again by Question 6,
\[
\mu\circ(\mu\times \mathrm{id}) \simeq \mu\circ(\mathrm{id}\times \mu).
\]
So \(\mu\) is associative up to homotopy. Now let \(X\) be any CW complex. For maps \(f,g:X\to Y\), define
\[
[f]+[g]:=[\,\mu\circ (f,g)\,],
\]
where
\[
(f,g):X\to Y\times Y,\qquad x\mapsto (f(x),g(x)).
\]
This gives an addition on \([X,Y]\).

We now show that under the canonical bijection
\[
[X,K(G,n)]\cong H^n(X;G),
\]
this addition agrees with the usual addition in cohomology. Let $\alpha,\beta\in H^n(X;G)$ correspond to $f,g:X\to K(G,n)$
By naturality of the representing-space correspondence, the product map $(f,g):X\to K(G,n)\times K(G,n)$
corresponds to the pair $(\alpha,\beta)\in H^n(X;G)\oplus H^n(X;G)$
The map $\mu:K(G,n)\times K(G,n)\to K(G,n)$
was chosen to represent the addition map $G\oplus G\to G$, so $\mu\circ(f,g)$ corresponds to the sum $\alpha+\beta\in H^n(X;G)$.
\end{solution}
\end{document}